% -----------------------------------------------------------------------------------------------------
\section{单圈计算：由算符定义得到 $P_{qq}$}\label{sec:pdf-oneloop}

本节在 Feynman 规范下直接由\cref{eq:pdf-fq}算出夸克中夸克分布的单圈值；Wilson 线显式出现，每个图的紫外极点可以单独读出。

\begin{proposition}[Feynman 规范下由算符定义计算 $P_{qq}$ \tagR]\label{prop:pdf-pqq-operator}
（$R_{\mathrm{PT}}$：在重整化微扰论中对红外正规化的部分子态严格。）
在单圈阶，取 Feynman 规范，以胶子质量 $m$ 作红外正规化、维数正规化作紫外正规化，靶为自旋与色平均的在壳无质量夸克 $p=(p^+,0^-,\bm 0_{\mathrm T})$，记
\begin{equation}\label{eq:pdf-Aeps}
A_\varepsilon\equiv\frac{\as}{2\pi}C_F\,\ee^{\varepsilon\gamma_E}\Gamma(\varepsilon)\Big(\frac{\mu^2}{m^2}\Big)^{\varepsilon}=\frac{\as}{2\pi}C_F\Big[\frac1\varepsilon+\ln\frac{\mu^2}{m^2}+\order{\varepsilon}\Big].
\end{equation}
(a) 实辐射部分为 $f_R(\xi)=A_\varepsilon\big[(1-\varepsilon)^2(1-\xi)\xi^{-\varepsilon}+\frac{2\xi^{1-\varepsilon}}{1-\xi}\big]$。
(b) Wilson 线的实图与虚图只在组合中构成分布 $A_\varepsilon[\frac{2\xi^{1-\varepsilon}}{1-\xi}]_+$（$\xi\to1$ 的端点奇异性不受维数正规化调控）；外线自能贡献 $-A_\varepsilon\frac{1-\varepsilon}{2-\varepsilon}\delta(1-\xi)$。
(c) 裸单圈分布对检验函数 $h$ 为
\begin{equation}\label{eq:pdf-bare1}
\langle f^{(1)}_{\mathrm{bare}},h\rangle=A_\varepsilon\Big\{\int_0^1\!\dd\xi\Big[(1-\varepsilon)^2(1-\xi)\xi^{-\varepsilon}h(\xi)+\frac{2\xi^{1-\varepsilon}}{1-\xi}\big(h(\xi)-h(1)\big)\Big]-\frac{1-\varepsilon}{2-\varepsilon}h(1)\Big\},
\end{equation}
夸克数对任意 $\varepsilon$ 精确守恒：$\int_0^1f^{(1)}_{\mathrm{bare}}=0$。
(d) 紫外极点逐图（单位 $\frac{\as}{2\pi}\frac{C_F}{\varepsilon}$）：梯形图 $(1-\xi)$，Wilson 线图 $[\frac{2\xi}{1-\xi}]_+$，自能 $-\frac12\delta(1-\xi)$；其和为 $P_{qq}/C_F$。
(e) $\MSbar$ 重整化后（依赖于红外正规化，不是物理量）
\begin{equation}\label{eq:pdf-ren1}
f^{(1)}_{q/q}(\xi;\mu)=\frac{\as}{2\pi}\Big\{P_{qq}(\xi)\ln\frac{\mu^2}{m^2}+C_F\Big[-(1-\xi)(2+\ln\xi)-\Big[\frac{2\xi\ln\xi}{1-\xi}\Big]_++\frac14\delta(1-\xi)\Big]\Big\}.
\end{equation}
(f) 在光锥规范中只有梯形图，它单独给出同一 $D$ 维核 $\hat P_{qq}$。
（本命题只涉及部分子态。）
% VF-PDF-3
\uses{def:pdf-quark,def:prelim-feynman-rules,prop:prelim-wilson-rules,def:prelim-msbar}
\end{proposition}

\begin{proof}
\emph{图。}用割图表示\cref{eq:pdf-cut}。$\order{\as}$ 的图见\cref{fig:pdf-oneloop}：(a) 梯形图 R1；(b) 实胶子一端在夸克线、另一端在割线另一侧的 Wilson 线上（R2 及镜像）；(c) 虚图：夸克线—Wilson 线顶点修正 V1 与外线自能 V2（及镜像）；(d) 两端都在 Wilson 线上的图 $\propto n_B^2=0$。
\emph{实辐射。}由\cref{prop:prelim-wilson-rules}，$W[\infty n_B,0]$ 吸收流入动量 $-l$（即发射出射胶子 $l$）的单胶子项为 $g\,t^an_B^\mu/(-n_B\cdot l-\ii0)$，故
\begin{equation}\label{eq:pdf-amp}
\langle l,a,\epsilon|W[\infty n_B,0]\psi(0)|p\rangle=-g\,t^a\epsilon^*_\mu\Big[\frac{\sla k\gamma^\mu}{k^2}+\frac{n_B^\mu}{n_B\cdot l+\ii0}\Big]u(p),\qquad k=p-l .
\end{equation}
相对符号由规范不变性确定：$\sla lu(p)=-\sla ku(p)$，故 $l_\mu[\cdots]u=(\sla k\sla l/k^2+1)u=0$；因此有质量胶子极化和中的 $l^\mu l^\nu/m^2$ 项不贡献，可用 $-g_{\mu\nu}$。运动学：$l=\big((1-\xi)p^+,\frac{\bm l_{\mathrm T}^2+m^2}{2(1-\xi)p^+},\bm l_{\mathrm T}\big)$，$n_B\cdot k=\xi p^+$，$k^2=-(\bm l_{\mathrm T}^2+\xi m^2)/(1-\xi)$。对靶自旋平均后，$D$ 维迹给出（$\Gamma^\mu$ 为\cref{eq:pdf-amp}方括号中的矩阵）
\[
\tfrac12\TrD[\sla p\,\bar\Gamma^\nu\gamma^+\Gamma^\mu](-g_{\mu\nu})g^2C_F=g^2C_Fp^+\Big[\underbrace{\frac{4(1-\varepsilon)(1-\xi)^2\bm l_{\mathrm T}^2}{(\bm l_{\mathrm T}^2+\xi m^2)^2}}_{\text{R1}}+\underbrace{\frac{8\xi}{\bm l_{\mathrm T}^2+\xi m^2}}_{\text{R2+镜像}}\Big],
\]
R1 来自 $\gamma^\mu\sla k\gamma^+\sla k\gamma_\mu=-(D-2)\sla k\gamma^+\sla k$，R2 来自 $\gamma^+\sla k\gamma^+=2k^+\gamma^+$。一粒子相空间中 $l^+$ 被 $\delta$ 函数固定，连同\cref{eq:pdf-cut}中的 $\frac12$，
\[
f_R(\xi)=\frac{g^2C_F}{(2\pi)^{D-1}}\int\dd^{D-2}\bm l_{\mathrm T}\Big[\frac{(1-\varepsilon)(1-\xi)\bm l_{\mathrm T}^2}{(\bm l_{\mathrm T}^2+\xi m^2)^2}+\frac{2\xi}{(1-\xi)(\bm l_{\mathrm T}^2+\xi m^2)}\Big].
\]
用 $\int\dd^{d}\bm l\,\frac{(\bm l^2)^a}{(\bm l^2+M^2)^b}=\pi^{d/2}\frac{\Gamma(a+d/2)\Gamma(b-a-d/2)}{\Gamma(d/2)\Gamma(b)}(M^2)^{a+d/2-b}$（$d=2-2\varepsilon$，$M^2=\xi m^2$）与 $\MSbar$ 中的 $\frac{g_0^2}{(2\pi)^{3-2\varepsilon}}\pi^{1-\varepsilon}=\frac{\as}{2\pi}\mu^{2\varepsilon}\ee^{\varepsilon\gamma_E}$ 即得 (a)。
\emph{虚 Wilson 线图 V1。}用 $\bar u\gamma^+(\sla p-\sla l)\gamma^+u=2(p^+-l^+)\bar u\gamma^+u$，
\[
\frac{V_1}{\text{树图}}=\ii g^2C_F\int\frac{\dd^Dl}{(2\pi)^D}\,\frac{2(p^+-l^+)}{(l^+-\ii0)(l^2-m^2+\ii0)((p-l)^2+\ii0)} .
\]
两个传播子在 $l^-$ 平面的极点 $\frac{\bm l_{\mathrm T}^2+m^2-\ii0}{2l^+}$ 与 $-\frac{\bm l_{\mathrm T}^2-\ii0}{2(p^+-l^+)}$ 只当 $0<l^+<p^+$ 时位于实轴两侧；闭合围道得 $\frac{\pi\ii\,\theta(0<l^+<p^+)}{p^+(\bm l_{\mathrm T}^2+xm^2)}$，$l^+=(1-x)p^+$。用 $\frac{2(p^+-l^+)}{l^+p^+}\dd l^+=\frac{2x}{1-x}\dd x$ 与 $\ii\cdot\frac{\pi\ii}{(2\pi)^D}=-\frac{1}{2(2\pi)^{D-1}}$，加上镜像，
\[
\frac{V_1+V_1^\dagger}{\text{树图}}=-\frac{g^2C_F}{(2\pi)^{D-1}}\int_0^1\!\dd x\!\int\!\dd^{D-2}\bm l_{\mathrm T}\frac{2x}{(1-x)(\bm l_{\mathrm T}^2+xm^2)}=-A_\varepsilon\int_0^1\!\dd x\,\frac{2x^{1-\varepsilon}}{1-x},
\]
它乘以 $\delta(1-\xi)$，被积函数恰为 R2 在 $\xi=x$ 处的相反数，所以 R2 与 V1 合成 $A_\varepsilon[2\xi^{1-\varepsilon}/(1-\xi)]_+$；端点 $x\to1$（$l^+\to0$）的发散只在这一组合中相消，必须在固定 $x=\xi$ 处先组合再积分。
\emph{外线自能 V2。}$\Sigma(p)=\sla p\,\Sigma_V(p^2)$，$\Sigma_V(0)=(2-D)\frac{g_0^2C_F}{(4\pi)^{D/2}}\Gamma(\varepsilon)\int_0^1\dd y\,(1-y)[(1-y)m^2]^{-\varepsilon}=-A_\varepsilon\frac{1-\varepsilon}{2-\varepsilon}$，其极点 $-\frac{\as}{4\pi}\frac{C_F}\varepsilon$ 即 Feynman 规范的 $\delta Z_2$；两条外线的 LSZ 留数合计贡献 $(R-1)\delta(1-\xi)=\Sigma_V(0)\delta(1-\xi)+\order{\as^2}$。这证明了 (b)，相加得 (c)。
\emph{夸克数、极点、有限部分。}$(1-\varepsilon)^2\int_0^1(1-\xi)\xi^{-\varepsilon}\dd\xi=\frac{1-\varepsilon}{2-\varepsilon}$，加号分布积分为零，故 $h=1$ 时\cref{eq:pdf-bare1}为零。由 $\ee^{\varepsilon\gamma_E}\Gamma(\varepsilon)=\frac1\varepsilon+\order\varepsilon$（无常数项），极点为 $\frac{\as}{2\pi}\frac{C_F}\varepsilon\{(1-\xi)+[\frac{2\xi}{1-\xi}]_+-\frac12\delta\}$，而用 $[\frac{2\xi}{1-\xi}]_+=\frac{2\xi}{(1-\xi)_+}+2\delta(1-\xi)$，花括号等于 $\frac{1+\xi^2}{(1-\xi)_+}+\frac32\delta=[\frac{1+\xi^2}{1-\xi}]_+$，即 (d)。减去 $\frac{\as}{2\pi}\frac1\varepsilon P_{qq}$ 后，$\varepsilon^0$ 项是 $\frac{\as}{2\pi}P_{qq}\ln\frac{\mu^2}{m^2}$ 加上花括号对 $\varepsilon$ 的一阶导数：$\frac{\dd}{\dd\varepsilon}[(1-\varepsilon)^2(1-\xi)\xi^{-\varepsilon}]_0=-(1-\xi)(2+\ln\xi)$，$\frac{\dd}{\dd\varepsilon}[2\xi^{1-\varepsilon}]_0=-2\xi\ln\xi$，$-\frac{\dd}{\dd\varepsilon}[\frac{1-\varepsilon}{2-\varepsilon}]_0=\frac14$，即 (e)。
\emph{光锥规范。}取 $d^{\mu\nu}(l)=-g^{\mu\nu}+(l^\mu n_B^\nu+n_B^\mu l^\nu)/(n_B\cdot l)$，$W=1$，只剩梯形图，其迹乘以 $\bm l_{\mathrm T}^2$（$m=0$）同样等于 $\frac{1+\xi^2}{1-\xi}-\varepsilon(1-\xi)$，即 (f)：Feynman 规范中由 Wilson 线承担的部分，在光锥规范中由极化和的 $1/(n_B\cdot l)$ 项承担。
\end{proof}

%%FIG:pdf-oneloop%%

\begin{example}[单圈夸克 PDF（逐图）]\label{ex:pdf-oneloop}
按\cref{fig:pdf-oneloop}逐图：R1 给出 $A_\varepsilon(1-\varepsilon)^2(1-\xi)\xi^{-\varepsilon}$，R2 与 V1 给出 $A_\varepsilon[\frac{2\xi^{1-\varepsilon}}{1-\xi}]_+$，V2 给出 $-A_\varepsilon\frac{1-\varepsilon}{2-\varepsilon}\delta(1-\xi)$，(d) 为零；$\delta$ 项 $\frac32=2-\frac12$ 中，$2$ 来自把 R2 与 V1 的组合改写为标准形式，$-\frac12$ 来自夸克场重整化，梯形图不含 $\delta$ 项。数值上，取 $\mu=m$、$\xi=\frac12$，\cref{eq:pdf-ren1}的正规部分为 $\frac{\as C_F}{2\pi}[-\frac12(2-\ln2)+2\ln2]=\frac{\as C_F}{2\pi}(\frac52\ln2-1)=0.7329\,\frac{\as C_F}{2\pi}$；$\ln(\mu^2/m^2)$ 每改变一个单位，该点的值改变 $\frac{\as}{2\pi}P_{qq}(\frac12)=\frac52\frac{\as C_F}{2\pi}$，正是 DGLAP 所要求的 $\mu$ 依赖。夸克数检验：$\int_0^1[-(1-\xi)(2+\ln\xi)]\dd\xi=-\frac14$，恰与 $\frac14\delta(1-\xi)$ 相消。有限部分依赖于红外正规化（纯维数正规化下为零，\cref{prop:pdf-partonic-dr}），只有紫外极点与 $\ln\mu^2$ 的系数与正规化无关。
\uses{prop:pdf-pqq-operator}
\end{example}

\begin{proposition}[纯维数正规化中的部分子分布 \tagR]\label{prop:pdf-partonic-dr}
（$R_{\mathrm{PT}}$：在重整化微扰论中对红外正规化的部分子态严格。）
设靶为无质量在壳部分子 $j$，维数正规化同时正规化紫外与红外发散。则裸部分子分布在所有阶都是 $f^{\mathrm{bare}}_{i/j}(\xi)=\delta_{ij}\delta(1-\xi)$；$\MSbar$ 重整化后
\begin{equation}\label{eq:pdf-partonicDR}
f_{i/j}(\xi;\mu)=Z_{ij}(\xi)=\delta_{ij}\delta(1-\xi)-\frac{\as}{2\pi}\frac1\varepsilon P_{ij}(\xi)+\order{\as^2},
\end{equation}
% VF-PDF-5
剩下的 $1/\varepsilon$ 是红外极点：$\int\frac{\dd^{2-2\varepsilon}\bm l_{\mathrm T}}{\bm l_{\mathrm T}^2}=\pi^{1-\varepsilon}\Gamma(1+\varepsilon)M^{-2\varepsilon}\big(\frac1{\varepsilon_{\mathrm{UV}}}-\frac1{\varepsilon_{\mathrm{IR}}}\big)=0$，$\MSbar$ 抵消项去掉 $1/\varepsilon_{\mathrm{UV}}$ 后留下 $-1/\varepsilon_{\mathrm{IR}}$。（本命题只涉及部分子态。）
\uses{def:pdf-quark,prop:pdf-invariance,lem:prelim-scaleless,ex:prelim-scaleless,prop:pdf-lo-kernels}
\end{proposition}

\begin{proof}
设 $F_L(\xi,p^+;\varepsilon)$ 是 $L\ge1$ 圈的裸贡献。它无量纲，唯一有量纲的变量是 $p^+$（$p^2=0$，无质量）；按量纲分析，把所有动量放大 $\lambda$ 倍时 $L$ 圈测度与 $g_0^{2L}\propto\mu_0^{2L\varepsilon}$ 给出 $F_L(\xi,\lambda p^+)=\lambda^{-2L\varepsilon}F_L(\xi,p^+)$；另一方面推动不变性（\cref{prop:pdf-invariance}(b)，对裸分布同样成立）给出 $F_L(\xi,\lambda p^+)=F_L(\xi,p^+)$。故 $(\lambda^{-2L\varepsilon}-1)F_L=0$ 对一切 $\lambda>0$ 成立，$\varepsilon\ne0$ 时 $F_L=0$——这正是\cref{lem:prelim-scaleless}的情形。所以 $f^{\mathrm{bare}}_{i/j}=\delta_{ij}\delta(1-\xi)$（\cref{prop:pdf-invariance}(c)），$f^{\mathrm{ren}}_{i/j}=\sum_kZ_{ik}\otimes\delta_{kj}\delta=Z_{ij}$，单圈的 $Z$ 由\cref{eq:pdf-Z}与\cref{prop:pdf-lo-kernels}给出\cref{eq:pdf-partonicDR}。单圈实辐射积分正是 $\int\dd^{2-2\varepsilon}\bm l_{\mathrm T}/\bm l_{\mathrm T}^2$：在任意 $M$ 处分开，$|\bm l_{\mathrm T}|>M$ 部分只在 $\varepsilon>0$ 时收敛，给出 $+1/\varepsilon_{\mathrm{UV}}$，$|\bm l_{\mathrm T}|<M$ 部分只在 $\varepsilon<0$ 时收敛，给出 $-1/\varepsilon_{\mathrm{IR}}$，各自延拓后相消（\cref{ex:prelim-scaleless}）。$\MSbar$ 抵消项与红外正规化无关，去掉 $1/\varepsilon_{\mathrm{UV}}$ 后剩下的 $-\frac{\as}{2\pi}\frac1{\varepsilon_{\mathrm{IR}}}P_{ij}$ 是部分子靶的共线发散。
\end{proof}

这正是把部分子截面变成系数函数的那一步：对部分子靶，因子化公式中的 PDF 是纯红外极点，所以系数函数等于“$\MSbar$ 减除了共线极点的部分子截面”（\cref{ch:disall}）；忘记\cref{eq:pdf-partonicDR}右边的极点是红外的，共线减除中就很容易弄错符号。

\begin{proposition}[定域算符的单圈反常维度（含混合）\tagR]\label{prop:pdf-local-anomalous}
（$R_{\mathrm{PT}}$：在重整化微扰论中对红外正规化的部分子态严格。）
在单圈阶，对在壳、物理极化的无质量部分子态 $j$ 与扭度 2 定域算符 $O_q^{(N)}=\bar\psi\sla n_B(\ii n_B\cdot D)^{N-1}\psi$（单一味道）、$O_g^{(N)}=\sum_jG^{+j}(\ii D^+)^{N-2}G^{+j}$，定义 $O_i^R=\sum_jZ_{ij}O_j$，$Z=\mathbb 1-\frac{\as}{2\pi}\frac1\varepsilon\Gamma$（于是 $\frac{\dd O^R}{\dd\ln\mu^2}=\frac{\as}{2\pi}\Gamma O^R$）。则 $\langle j|O_i|j\rangle$ 的紫外极点给出 $\Gamma_{qq}=\gamma_{qq}(N)$，$\Gamma_{q_kg}=2\gamma_{qg}(N)$，$\Gamma_{gq}=\gamma_{gq}(N)$，$\Gamma_{gg}=\gamma_{gg}(N)$（偶数 $N$；奇数 $N$ 只有非单态 $\Gamma_{qq}=\gamma_{qq}(N)$）。Feynman 规范中逐图的值见\cref{tab:pdf-localgraphs}；单个图依赖于规范参数 $\xi_g$，只有包括 $\delta Z_2$、$\delta Z_3$ 的总和与 $\xi_g$ 无关。$N=2$ 时 $\Gamma$ 的两列之和为零：$T^{++}=\sum_kO^{(2)}_{q_k}+O^{(2)}_g$ 不被重整化。（本命题只涉及部分子态与算符。）
% VF-PDF-14, VF-PDF-10
\uses{def:dis-twist,thm:prelim-composite-renorm,prop:pdf-lo-kernels}
\end{proposition}

\begin{proof}
\emph{重整化条件。}树图 $\langle j|O_i|j\rangle^{(0)}=\delta_{ij}\cdot2(n_B\cdot p)^N$（偶数 $N$）。由\cref{thm:prelim-composite-renorm}，对物理外态扭度 2 算符只在 $\{O_q,O_g\}$ 内混合，于是 LSZ 振幅（含外线留数）的有限性要求 $\mathrm{UV}[\Gamma^{(1)}_{\mathrm{1PI}}+\delta Z_j\langle j|O_i|j\rangle^{(0)}]/(2(n_B\cdot p)^N)=\frac{\as}{2\pi}\frac1\varepsilon\Gamma_{ij}$，$\delta Z_q=\delta Z_2$，$\delta Z_g=\delta Z_3$。
\emph{紫外极点的提取。}单圈紫外发散定域、与红外正规化无关。对无质量在壳运动学，Feynman 参数化并平移后分母是 $q^2$ 的纯幂，只有分子中关于回路动量 $q$ 为二次（三个传播子）或零次（两个传播子）的项是对数发散的；用 $\langle q^\mu q^\nu\rangle=q^2g^{\mu\nu}/D$ 与 $\int\frac{\dd^Dq}{(2\pi)^D}\frac{q^2}{(q^2)^3}\big|_{\mathrm{UV}}=\int\frac{\dd^Dq}{(2\pi)^D}\frac{1}{(q^2)^2}\big|_{\mathrm{UV}}=\frac{\ii}{16\pi^2\varepsilon}$。下面记 $n=n_B$。
\emph{夸克部分。}顶点图 A1 $=-\ii g^2C_F\int\frac12\TrD[\sla p\gamma^\mu\sla k\sla n\sla k\gamma_\mu](n\cdot k)^{N-1}/((k^2)^2(p-k)^2)$；用 $\gamma^\mu\sla k\sla n\sla k\gamma_\mu=-(D-2)\sla k\sla n\sla k$、权重 $2(1-y)$、$k=q+yp$，分子的二次部分为 $2y^{N-1}(n\cdot p)^Nq^2$，极点为 $\frac{\as}{4\pi}\frac{C_F}\varepsilon(n\cdot p)^N\int_0^12(1-y)2y^{N-1}\dd y$，相对树图为 $\frac{\as}{2\pi}\frac{C_F}\varepsilon\frac1{N(N+1)}$。算符—胶子图 A2 $=$ A3 $=\ii g^2C_F\int\frac12\TrD[\sla p\sla n\sla k\sla n]\sum_{j=0}^{N-2}(n\cdot p)^j(n\cdot k)^{N-2-j}/(l^2k^2)$，$k=p-l$；由 $\sla n\sla k\sla n=2(n\cdot k)\sla n$ 与 $l=q+xp$，只剩与 $q$ 无关的项，极点为 $-\frac{\as}{4\pi}\frac{C_F}\varepsilon4(n\cdot p)^N\sum_{m=2}^N\frac1m$，即每个图 $-\frac{\as}{2\pi}\frac{C_F}{\varepsilon}(S_1(N)-1)$。加上 Feynman 规范的 $\delta Z_2=-\frac{\as}{4\pi}\frac{C_F}{\varepsilon}$：$\Gamma_{qq}=C_F[\frac{1}{N(N+1)}-2(S_1(N)-1)-\frac12]=\gamma_{qq}(N)$。
\emph{混合与胶子部分。}$\langle q|O_g|q\rangle$ 的单个图 B 给出 $\gamma_{gq}$；$\langle g|O_q|g\rangle$ 中夸克圈的两种走向 C1、C2 各给出 $\gamma_{qg}$，且由 C 宇称 $C_2=(-1)^NC_1$，故偶数 $N$ 时 $\Gamma_{q_kg}=2\gamma_{qg}$，奇数 $N$ 时为零（$N=1$ 即 Furry 定理）。D1（算符的双胶子顶点与两个三胶子顶点）、D2 $=$ D3（算符的三胶子顶点）与 $\delta Z_3=\frac{\as}{4\pi}\frac1\varepsilon[(\frac{13}6-\frac{\xi_g}2)C_A-\frac43T_Fn_f]$ 之和为 $\gamma_{gg}$；$n_f$ 部分全部来自 $\delta Z_3$。由算符双胶子项构成的蝌蚪图与“算符双胶子顶点 $+$ 四胶子顶点”图没有紫外极点。
\emph{规范无关性与 $N=2$。}表中 $\xi_g$ 依赖之和为零（夸克：$\frac{\xi_g-1}2-\frac{\xi_g-1}2=0$；胶子 $C_A$ 部分：$\frac{\xi_g-1}2-2\cdot\frac{\xi_g-1}8-\frac{\xi_g-1}4=0$）。$N=2$ 时 $\Gamma=\begin{pmatrix}-\frac43C_F&\frac23n_fT_F\\\frac43C_F&-\frac23n_fT_F\end{pmatrix}$（基 $\sum_kO_{q_k}$、$O_g$），两列之和为零。
% (graph values cross-checked for N=1..14 with explicit 4x4 Dirac matrices and with D-dim traces; writing/work/W09/v9_local_ops.py)
\end{proof}

\begin{table}[htbp]
\centering
\caption{扭度 2 定域算符单圈矩阵元的逐图紫外极点（单位 $\frac{\as}{2\pi}\frac1\varepsilon\times$ 色因子，相对于树图 $2(n_B\cdot p)^N$），以及一般协变规范中相对于 Feynman 规范（$\xi_g=1$）的改变。}\label{tab:pdf-localgraphs}
\small
\renewcommand{\arraystretch}{1.2}
\begin{tabular}{@{}lll@{}}
\toprule
图 & Feynman 规范 & 一般协变规范的改变\\
\midrule
A1 顶角（$C_F$） & $\frac{1}{N(N+1)}$ & $+\frac{\xi_g-1}{2}$\\
A2 $=$ A3 算符—胶子（$C_F$） & $-(S_1(N)-1)$ & 无\\
$\delta Z_2$（$C_F$） & $-\frac12$ & $-\frac{\xi_g-1}{2}$\\
B（$C_F$） & $\frac{N^2+N+2}{(N-1)N(N+1)}$ & 无\\
C1 $=$ C2（$T_F$，每味） & $\frac{N^2+N+2}{N(N+1)(N+2)}$ & 无（无内胶子线）\\
D1 双胶子算符顶角（$C_A$） & $\frac{2}{N(N-1)}+\frac{2}{(N+1)(N+2)}+\frac{1}{N+1}$ & $+\frac{\xi_g-1}{2}$\\
D2 $=$ D3 三胶子算符顶角（$C_A$） & $-(S_1(N)-1)-\frac{N+2}{2(N+1)}$ & $-\frac{\xi_g-1}{8}$\\
$\delta Z_3$ & $\frac56C_A-\frac23T_Fn_f$ & $-\frac{\xi_g-1}{4}C_A$\\
\bottomrule
\end{tabular}
\end{table}

\begin{example}[精确到 $\varepsilon$ 的单圈矩]\label{ex:pdf-moments-exact}
取\cref{eq:pdf-bare1}的 $n$ 阶矩，并与同样以胶子质量正规化的定域算符单圈矩阵元比较，PDF 图与定域算符图逐一对应且对 $\varepsilon$ 精确相等（单位 $A_\varepsilon$）：梯形图 R1 $\leftrightarrow$ 顶角图 A1，$(1-\varepsilon)^2\int_0^1\xi^{n-1}(1-\xi)\xi^{-\varepsilon}=\frac{(1-\varepsilon)^2}{(n-\varepsilon)(n+1-\varepsilon)}$；Wilson 线图 R2$+$V1 $\leftrightarrow$ 算符—胶子图 A2$+$A3，$\int_0^1\frac{2\xi^{1-\varepsilon}(\xi^{n-1}-1)}{1-\xi}=-2\sum_{i=1}^{n-1}\frac{1}{i+1-\varepsilon}$；自能 $\leftrightarrow$ 自能，$-\frac{1-\varepsilon}{2-\varepsilon}$。
% VF-PDF-13
总和：$n=1$ 时为 $0$（夸克数守恒对任意 $\varepsilon$ 成立）；$n=2$ 时 $\frac{(1-\varepsilon)^2-2(3-\varepsilon)-(1-\varepsilon)(3-\varepsilon)}{(2-\varepsilon)(3-\varepsilon)}=-\frac{4}{3-\varepsilon}$，极点 $-\frac43\frac{\as}{2\pi}\frac{C_F}\varepsilon=\frac{\as}{2\pi}\frac{\gamma_{qq}(2)}\varepsilon$；$n=3$ 时 $-\frac{2(5-2\varepsilon)^2}{(4-\varepsilon)(3-\varepsilon)(2-\varepsilon)}$，极点系数 $-\frac{25}{12}=\gamma_{qq}(3)/C_F$。一般地总和在 $\varepsilon\to0$ 的值 $\frac{1}{n(n+1)}-2(S_1(n)-1)-\frac12=\gamma_{qq}(n)/C_F$，于是 PDF 抵消项 $Z(x)=\delta(1-x)-\frac{\as}{2\pi}\frac1\varepsilon P_{qq}$ 的矩恰是定域算符的 $Z_O(n)$——\cref{thm:pdf-moments-local}第 (2) 步在单圈的显式实例。
\uses{thm:pdf-moments-local,prop:pdf-pqq-operator}
\end{example}

\begin{example}[胶子中的夸克]\label{ex:pdf-quark-in-gluon}
单圈 $f_{q/g}$ 来自 $g\to q\bar q$，$D$ 维核 $\hat P_{qg}=T_F[1-\frac{2z(1-z)}{1-\varepsilon}]$（\cref{prop:pdf-lo-kernels}）。无质量夸克的纯维数正规化中 $\MSbar$ 分布是纯红外极点 $-\frac{\as}{2\pi}\frac1\varepsilon P_{qg}$。为看清极化平均约定的作用，给夸克质量 $m$，$A^+=0$ 规范的直接计算给出横向被积函数 $\frac{g^2T_F}{(2\pi)^{D-1}}\frac{(\hat P_{qg}/T_F)\bm l_{\mathrm T}^2+m^2}{(\bm l_{\mathrm T}^2+m^2)^2}$（$m^2$ 项来自分子中的夸克质量），于是
\[
f^{(1)}_{Q/g}(z)=\frac{\as}{2\pi}\ee^{\varepsilon\gamma_E}\Big(\frac{\mu^2}{m^2}\Big)^{\varepsilon}\big[(1-\varepsilon)\hat P_{qg}\,\Gamma(\varepsilon)+T_F\Gamma(1+\varepsilon)\big]
=\frac{\as}{2\pi}P_{qg}(z)\,\ee^{\varepsilon\gamma_E}\Gamma(\varepsilon)\Big(\frac{\mu^2}{m^2}\Big)^{\varepsilon},
\]
因为 $(1-\varepsilon)\hat P_{qg}=P_{qg}-\varepsilon T_F$ 而 $\varepsilon\Gamma(\varepsilon)=\Gamma(1+\varepsilon)$：质量项恰好抵消了 $D$ 维核的 $\order{\varepsilon}$ 部分。$\MSbar$ 重整化后 $f^{(1)}_{Q/g}=\frac{\as}{2\pi}P_{qg}(z)\ln\frac{\mu^2}{m^2}$，没有常数项；若把胶子极化平均改为 $\frac12$，结果乘以 $\frac{D-2}{2}=1-\varepsilon$，$\MSbar$ 分布变为 $\frac{\as}{2\pi}P_{qg}[\ln\frac{\mu^2}{m^2}-1]$：极化平均约定是 $\MSbar$ 方案的一部分（\cref{def:pdf-gluon}）。$\mu<m$ 时这个 $\MSbar$ 重夸克分布是负的：$\MSbar$ 分布不必为正，微扰论中就有初等的例子（\cref{conj:pdf-positivity}）。
\uses{def:pdf-gluon,prop:pdf-lo-kernels}
\end{example}
